\chapter{Polynomial Congruences}\label{chapter:polynomial-congruences}

\section{Polynomial Congruences of Composite Moduli}

In the following theorem we generalize \zcrefnamed{theorem:chinese-remainder-theorem} to polynomial congruences.

\begin{theorem}
  Consider an integer polynomial $f$ and consider two relatively prime integers $m$ and $n$. Then any solution of the system
  \begin{align*}
    f(x) & \equiv 0 \pmod{m} \\
    f(x) & \equiv 0 \pmod{n} \\  
  \intertext{of polynomial congruences corresponds uniquely to a solution of the polynomial congruence}
    f(x) & \equiv 0 \pmod{mn}.
  \end{align*}
\end{theorem}

\begin{proof}
  First note that if $f(x)\equiv 0\pmod{mn}$, then it follows that $f(x)\equiv 0\pmod{m}$ and $f(x)\equiv 0\pmod{n}$, simply because divisibility is transitive.
  Consequently, if $y$ is the unique nonnegative integer strictly below $m$ such that $x\equiv y\pmod{m}$, then $f(y)\equiv \pmod{m}$ is also a solution, and similar for the modulus $n$. Thus if we write $N_{mn}(f)$ for the set of roots of $f$ modulo $mn$, and $N_m(f)$ and $N_n(f)$ for the sets of roots of $f$ modulo $m$ and $n$, respectively, then we define a map
  \begin{equation*}
    \pi : N_{mn}(f)\to N_m(f)\times N_n(f)
  \end{equation*}
  by $x\mapsto (x\pmod{m},x\pmod{n})$. We claim that the map $\pi$ is a bijection.

  To see that the map $\pi$ is a bijection, consider $0\leq y<m$ and $0\leq z<n$ such that
  \begin{align*}
    f(y) & \equiv 0 \pmod{m} \\
    f(z) & \equiv 0 \pmod{n}.
  \end{align*}
  By the assumption that $\gcd(m,n)=1$, we obtain from \zcrefnamed{theorem:chinese-remainder-theorem} there a unique $0\leq x < mn$ such that
  \begin{align*}
    x \equiv y \pmod {m} \\
    x \equiv z \pmod {n}.
  \end{align*}
  Consequently, we have that $f(x)\equiv 0\pmod{m}$ and $f(x)\equiv 0\pmod{n}$.
  Applying the Chinese Remainder Theorem once more to the number $f(x)$, using that $0$ is the only integer $u$ modulo $mn$ such that $u\equiv 0\pmod{m}$ and $u\equiv 0\pmod{n}$, we find that $f(x)\equiv 0\pmod{mn}$.
\end{proof}

\begin{corollary}
  Consider $n=p_1^{k_1}\cdots p_m^{k_m}$, where all the primes $p_i$ are distinct. A polynomial congruence
  \begin{equation*}
    f(x)\equiv 0 \pmod{n}
  \end{equation*}
  is solvable if and only if the polynomial congruence
  \begin{equation*}
    f(x)\equiv 0 \pmod{p_i^{k_i}}
  \end{equation*}
  is solvable for each $i$.
\end{corollary}

\begin{proof}
  By induction on the number of prime factors of $n$.
\end{proof}

\section{Derivatives of Polynomials}

The derivative of a function $f$ at a point $x$, if it exists, measures the rate of change of $f$ in close vicinity of $x$. It is formally defined as the limit
\begin{equation*}
  f'(x):=\lim_{h\to 0}\frac{f(x+h)-f(x)}{h}
\end{equation*}
Intuitively, such a limit exists if we can determine a unique value $y=f'(x)$ so that the quanitity $(f(x+h)-f(x))/h$ gets arbitrarily close to $y$ as $h$ gets arbitrarily close to $0$. When such a limit exists for every $x$, we say that $f$ is \emph{differentiable}. Many of the most common functions are differentiable, except at obvious points of discontinuity as seen, for instance, in step functions, or functions with sudden change of direction such as the absolute value function. In this short section we will prove that all polynomials are differentiable. There is an important theorem in number theory that require knowledge of the derivatives of polynomials: In our proof of \zcrefnamed{theorem:hensels-lemma} we will make use of \zcrefnamed{theorem:taylor-expansion-theorem} proven below.

The following theorem is not only useful when $h$ approaches $0$. For example, in the case where $h$ is a prime $p$ it tells us that $f(x+p)-f(x)$ is always a multiple of $p$.

\begin{theorem}
  Consider the polynomial
  \begin{equation*}
    f(x)=a_nx^n+a_{n-1}x^{n-1}+\cdots+a_1x+a_0.
  \end{equation*}
  Then the polynomial $f(x+h)-f(x)$ is divisible by $h$, with quotient
  \begin{equation*}
    \frac{f(x+h)-f(x)}{h} = \sum_{k=1}^n\sum_{i=0}^{k-1}a_kx^i(x+h)^{k-1-i}.
  \end{equation*}
\end{theorem}

\begin{proof}
  By the formula for the difference of $k$th powers, which was stated in \zcref{exercise:difference-of-powers}, we can write
  \begin{equation*}
    (x+h)^k-x^k= h\sum_{i=0}^{k-1}x^i(x+h)^{k-1-i}
  \end{equation*}
  In particular, since each of these differences of $k$th powers is divisible by $h$, we see that the polynomial $f(x+h)-f(x)$ is divisible by $h$. It follows that
  \begin{equation*}
    f(x+h)-f(x) = \sum_{k=0}^na_k((x+h)^k-x^k) = h\sum_{k=1}^na_k\sum_{i=0}^{k-1}x^i(x+h)^{k-1-i}. \qedhere
  \end{equation*}
\end{proof}

\begin{theorem}
  Consider the polynomial $f(x)=a_nx^n+\cdots+a_0$. Then $f$ is differentiable at every $x$, and the derivative can be computed by the formula
  \begin{equation*}
    f'(x) = \sum_{k=1}^n ka_kx^{k-1}.
  \end{equation*}
\end{theorem}

\begin{proof}
  Since polynomials are continuous functions, we can compute the limit as $h\to 0$ simply by evaluating the quotient of $f(x+h)-f(x)$ divided by $h$ at $0$:
  \begin{equation*}
    f'(x)=\sum_{k=1}^n\sum_{i=0}^{k-1}a_kx^i(x+0)^{k-1-i}=\sum_{k=1}^n\sum_{i=0}^{k-1}a_kx^{k-1}=\sum_{k=1}^n ka_kx^{k-1}
  \end{equation*}
  In the last equality we used that the summands in the inner summation do not depend on the summation index $i$.
\end{proof}

The derivative of any polynomial $f(x)$ is again a polynomial. An immediate consequence of the previous theorem is that if $h(x)=af(x)+bg(x)$ is a linear combination of the polynomials $f(x)$ and $g(x)$, then the derivative of $h$ can be written as the linear combination
\begin{equation*}
  h'(x)=af'(x)+bg'(x).
\end{equation*}
It also follows that, since the derivative is again a polynomial, it is again differentiable. Thus we have iterated derivatives of polynomials. We will write $\frac{d^k}{dx^k}f(x)$ or $f^{(k)}(x)$ for the $k$th deriviative of $f(x)$. A straightforward induction argument shows that
\begin{equation*}
  \frac{d^k}{dx^k}x^n = n(n-1)\cdots (n-k+1)x^{n-k}.
\end{equation*}
We can phrase this identity equivalently as:
\begin{equation*}
  \frac{1}{k!}\frac{d^k}{dx^k}x^n = \binom{n}{k}x^{n-k}.
\end{equation*}
This leads to the following theorem by Brook Taylor, who stated the result more generally in terms of \emph{fluent quantities} in his work \emph{Methodus incrementorum directa et inversa} \cite{taylor1715}. Although Taylor was the first to prove the series expansion that now bears his name, the method was discovered first without proof by James Gregory, who shared his insight in a letter to John Collins in 1671 \cite{gregory1671}. More on the history of Taylor series can be found in \cite{feigenbaum1985}.

\begin{theorem}[Taylor Expansion Theorem for Polynomials]\label{theorem:taylor-expansion-theorem}
  Consider the polynomial
  \begin{equation*}
    f(x)=a_nx^n+a_{n-1}x^{n-1}+\cdots+a_1x+a_0.
  \end{equation*}
  Then we have for each $y$, the identity
  \begin{equation*}
    f(x+y)=\sum_{k=0}^n \frac{f^{(k)}(x)}{k!}y^{k}.
  \end{equation*}
  In particular, there is a polynomial $g(x,y)$ in the variables $x$ and $y$, such that
  \begin{equation*}
    f(x+y)=f(x)+yf'(x)+y^2g(x,y).
  \end{equation*}
\end{theorem}

\begin{proof}
  We will first prove the theorem for $f(x):=x^n$. By \zcrefnamed{theorem:binomial-theorem} and our earlier observation that $\frac{1}{k!}\frac{d^k}{dx^k}x^n=\binom{n}{k}x^{n-k}$, it follows that
  \begin{equation*}
    (x+y)^n = \sum_{k=0}^{n}\binom{n}{k} x^{n-k}y^{k}=\sum_{k=0}^n\frac{1}{k!}\frac{d^kx^n}{dx^k} y^k.
  \end{equation*}
  This shows that the Taylor expansion formula holds for $f(x)=x^n$. The general claim follows, since every polynomial is a linear combination of polynomials of the form $x^i$ and differentiation preserves linear combinations. Finally, observe that
  \begin{equation*}
    f(x+y)=f(x)+yf'(x) + y^2\sum_{k=2}^n\frac{f^{(k)}(x)}{k!}y^{k-2}.
  \end{equation*}
  This completes the proof of the final claim.
\end{proof}

\section{Hensel's Lemma}

\begin{theorem}[Hensel's Lemma]\label{theorem:hensels-lemma}
  Consider a prime number $p$ and a polynomial $f(x)$ with integer coefficients. Furthermore, consider an integer $a$ such that
  \begin{equation*}
    f(a)\equiv 0\pmod{p}\qquad\text{and}\qquad f'(a)\nequiv 0\pmod{p},
  \end{equation*}
  where $f'(x)$ is the derivative of $f$. Then there exists for every $k\geq 1$ a unique integer $a_k$ modulo $p^k$ with $a_1\equiv a\pmod{p}$ such that the congruences
  \begin{equation*}
    a_{k+1}\equiv a_k\pmod{p^k},\qquad\text{and}\qquad f(a_k)\equiv 0\pmod{p^k}
  \end{equation*}
  hold for all $k$.
\end{theorem}

\begin{proof}
  We prove the claim by induction. That is, we prove by induction that for each $k\geq 1$ and every integer $a_k\equiv a\pmod{p}$ satisfying
  \begin{equation*}
    f(a_k)\equiv 0\pmod{p^k} \qquad\text{and}\qquad f'(a_k)\nequiv 0\pmod{p^k},
  \end{equation*}
  there is a unique integer $a_{k+1}\pmod{p^{k+1}}$ satisfying
  \begin{equation*}
    f(a_{k+1})\equiv 0\pmod{p^{k+1}}\qquad\text{and}\qquad f'(a_{k+1})\nequiv 0\pmod{p^{k+1}}.
  \end{equation*}
  Our goal is to obtain an integer $a_{k+1}=a_k+tp^k$ for some $0\leq t<p$. Observe that by \zcrefnamed{theorem:taylor-expansion-theorem}, we have the congruence
  \begin{equation*}
    f(a_k+tp^k)\equiv f(a_k)+ tp^k f'(a_k) \pmod{p^{k+1}}.
  \end{equation*}
  Since $f(a_k)\equiv 0\pmod{p^k}$ and $f'(a_k)\nequiv 0\pmod{p^k}$, it follows that there is a unique integer $t_0$ such that
  \begin{equation*}
    f(a_k)+ t_0p^k f'(a_k) \equiv 0 \pmod{p^{k+1}}.
  \end{equation*}
  We define $a_{k+1}:=a_k+t_0p^k$. Then we have $f(a_{k+1})\equiv 0\pmod{p^{k+1}}$ by definition. Furthermore, we have $a_{k+1}\equiv a_k\equiv a\pmod{p}$. This implies that $f'(a_{k+1})\equiv f'(a)\nequiv 0\pmod{p}$. Therefore it follows that $f'(a_{k+1})\nequiv 0\pmod{p^{k+1}}$, completing the proof.
\end{proof}

\section{The Elementary Symmetric Polynomials}

Using the fact that equivalent reduced polynomials must be congruent, we obtain the following factorization of the polynomial $x^{p-1}-1$ in $\mathbb{Z}/p\mathbb{Z}$, which has roots at the numbers $1$, \ldots, $p-1$.

\begin{theorem}\label{theorem:factorization-fermats-polynomial}
  Consider a prime $p$. Then we have
  \begin{equation*}
    x^{p-1}-1 \equiv (x-1)\cdots (x-p+1)\pmod{p}
  \end{equation*}
  for every $x$. Hence their coefficients are congruent.
\end{theorem}

\begin{proof}
  Both polynomials are equivalent and reduced, so they must be congruent.
\end{proof}

The previous theorem provides another way of proving \zcrefnamed{theorem:wilsons-theorem}. This proof is due to Lagrange, who was the first to publish a proof of Wilson's theorem.

\begin{proof}[Proof of \zcrefnamed{theorem:wilsons-theorem}]
  The integer $(p-1)!$ is the constant coefficient of the polynomial $(x-1)\cdots(x-p+1)$, which is congruent to the polynomial $x^{p-1}-1$ with constant coefficient $-1$.
\end{proof}

Using the elementary symmetric polynomials, we find another way of stating \zcref{theorem:factorization-fermats-polynomial}:

\begin{theorem}
  For each prime $p$ and each $1\leq k < p-1$ we have
  \begin{equation*}
    e_k(1,\ldots,p-1) \equiv 0 \pmod{p}.
  \end{equation*}
\end{theorem}

\section{Reduced Polynomials Modulo a Prime}

\begin{definition}
  Two polynomials $f(x)$ and $g(x)$ are said to be \emph{congruent} modulo a prime $p$ if their coefficients are congruent modulo $p$. 
\end{definition}

\begin{definition}
  Two polynomials $f(x)$ and $g(x)$ are said to be \emph{equivalent} modulo a prime $p$ if for each integer $x$ the polynomial congruence
  \begin{equation*}
    f(x)\equiv g(x)\pmod{p}
  \end{equation*}
  holds.
\end{definition}

\begin{definition}
  A polynomial $f(x)$ is said to be \emph{reduced} modulo $p$ if it is of degree less than $p$.
\end{definition}

\begin{theorem}
  Any two equivalent reduced polynomials modulo $p$ are congruent.
\end{theorem}

\begin{proof}
  Consider two equivalent reduced polynomials $f(x)$ and $g(x)$. Then $f(x)-g(x)$ is a reduced polynomial equivalent to $0$. However, if $f(x)-g(x)$ were nonzero then its number of roots would be strictly less than $p$. Since the number of roots is $p$, it follows that $f(x)-g(x)$ is congruent to the zero polynomial.
\end{proof}

\begin{theorem}
  For every polynomial $f(x)$ there is a unique reduced polynomial $g(x)$ equivalent to $f(x)$.
\end{theorem}

%%% Local Variables:
%%% mode: latex
%%% TeX-master: "elementary-number-theory"
%%% TeX-engine: luatex
%%% End:
